%%
%% fall-lecture10.tex
%% 
%% Made by Alex Nelson <pqnelson@gmail.com>
%% Login   <alex@lisp>
%% 
%% Started on  2026-10-20T11:53:55-0700
%% Last update 2026-10-20T11:53:55-0700
%% 

\lecture{}

\begin{node}[Recap---where are we?]
We saw (\S\ref{ex:w-ts-is-one}) that $\forall n\geq0\ldotp w(\T\sphere{n})=1$.
We are interested in $w(\T\RP^{n})$.
We proved (\S\ref{thm:lec09:structure-theorem}) $\T\RP^{n}\iso\hom(\gamma^{1}_{n},(\gamma^{1}_{n})^{\perp})$
but do not know how to compute $w\bigl(\hom(\gamma^{1}_{n},(\gamma^{1}_{n})^{\perp})\bigr)$.
\end{node}

\begin{theorem}
Let us denote $\varepsilon^{1}$ to be hte trivial line bundle over $\RP^{n}$.
Then
\begin{equation}
\T\RP^{n}\oplus\varepsilon^{1}\iso\underbrace{\gamma^{1}_{n}\oplus\dots\oplus\gamma^{1}_{n}}_{(n+1)\text{~times}}
\end{equation}
\end{theorem}

\begin{proof}
We have
\begin{equation}
\T\RP^{n}\oplus\varepsilon^{1}=\hom(\gamma^{1}_{n},(\gamma^{1}_{n})^{\perp})\oplus\hom(\gamma^{1}_{n},\gamma^{1}_{n}).
\end{equation}
We will use three facts (which we will not prove) in this proof:

\textsc{Fact 1}: $\RP^{n}\times\RR=\varepsilon^{1}\iso\hom(\gamma^{1}_{n},\gamma^{1}_{n})$
are isomorphic as bundles, with the isomorphism given by
$([x],t)\mapsto(x\mapsto tx)$.

\textsc{Fact 2}: $\hom(\xi,\eta)\oplus\hom(\xi,\eta')\iso\hom(\xi,\eta\oplus\eta')$
as bundles.

\textsc{Fact 3}: $(\gamma^{1}_{n})^{*}\iso\gamma^{1}_{n}$.

Then
\begin{subequations}
  \begin{align}
\T\RP^{n}\oplus\varepsilon^{1}
&=\hom(\gamma^{1}_{n},(\gamma^{1}_{n})^{\perp})\oplus\hom(\gamma^{1}_{n},\gamma^{1}_{n})\\
&=\hom(\gamma^{1}_{n},(\gamma^{1}_{n})^{\perp}\oplus\gamma^{1}_{n})\quad\mbox{by \textsc{Fact 2}}\\
&=\hom(\gamma^{1}_{n},\RP^{n}\times\RR^{n+1})\\
&=\hom(\gamma^{1}_{n},\bigoplus^{n+1}_{j=1}\varepsilon^{1})\\
&=\bigoplus^{n+1}_{j=1}\hom(\gamma^{1}_{n},\varepsilon^{1})\quad\mbox{by \textsc{Fact~2}}\\
&=\bigoplus^{n+1}_{j=1}(\gamma^{1}_{n})^{*}\\
&=\bigoplus^{n+1}_{j=1}\gamma^{1}_{n}\quad\mbox{using \textsc{Fact 3} and Riesz representation theorem}.
  \end{align}
\end{subequations}
We want to use Riesz representation theorem, but this requires a
Euclidean metric.

Then
\begin{subequations}
  \begin{align}
w(\T\RP^{n})
&=w(\T\RP^{n}\oplus\varepsilon^{1})\quad\mbox{by Axiom~\ref{axiom3}}\\
&=w(\gamma^{1}_{n})^{n+1}\quad\mbox{by Axiom~\ref{axiom3}}\\
&=(1+a)^{n+1}\quad\mbox{by Axiom~\ref{axiom4}}\\
&=1+\binom{n+1}{1}a+\binom{n+1}{2}a^{2}+\dots+\binom{n+1}{n}a^{n},
  \end{align}
\end{subequations}
since $a^{n+1}=0$.
\end{proof}

\begin{corollary}[Stiefel]
$w(\T\RP^{n})=0$ if and only if $n+1$ is a power of $2$.
\end{corollary}

\begin{proof}
\backwardproof\ Assume $n+1=2^{\ell}$, then
\begin{equation}
(1+a)^{n+1}=(1+a)^{2^{\ell}}=1+a^{n+1}=1.
\end{equation}

\forwardproof\ Proof by contrapositive: assume $n+1$ is not a power of
2. Then $n+1=m2^{k}$ where $m>1$ odd. Then expand
\begin{subequations}
  \begin{align}
w(\T\RP^{n}) &= (1+a)^{m2^{k}}=(1+a^{2^{k}})^{m}\\
&=1+ma^{2^{k}} + \bigOh(a^{2^{k}+1}).
  \end{align}
\end{subequations}
Since $2^{k}<n+1$, $a^{2^{k}}\neq0$, so $w(\T\RP^{n})\neq1$.
\end{proof}

\begin{remark}
If we assume that $\RP^{n}$ is parallelizable [so $\T\RP^{n}$ is
  trivial], then $w(\T\RP^{n})=1$. By our corollary, $n+1$ must be a
power of 2. But the converse is not true: for example
$w(\T\RP^{15})=1$ since $\T\RP^{15}$ is a pullback bundle of
$\T\sphere{15}$, but $\RP^{15}$ is 
not parallelizable. How do we know? Well, we use the quotient map $f\colon\sphere{15}\to\RP^{15}$
then $\T\sphere{15}\iso f^{*}\T\RP^{15}$ but $\T\sphere{15}$ is not
trivial.

The only parallelizable real projective spaces are $\RP^{1}$,
$\RP^{3}$, and $\RP^{7}$.
\end{remark}

\begin{remark}
Recall if we have a smooth manifold $M$, then $w_{1}(\T M)=0$ iff $M$
is orientable. [Homework: prove this.]

This can be generalized to bundles: $\xi$ is orientable if and only if
$w_{1}(\xi)=0$.

Now, $w_{1}(\T\RP^{n})=(n+1)a$, we can infer immediately $\RP^{n}$ is
orientable if and only if $n$ is odd.
\end{remark}

\subsection{Stiefel--Whitney Numbers}

\begin{recall}
Let $M$ be a closed smooth manifold and $\dim(M)=n$.
Then $H_{n}(M;\ZZ_{2})\iso\ZZ_{2}\oplus\dots\oplus\ZZ_{2}$ where there
are $|\pi_{0}(M)|$ summands (i.e., as many summands as there are
connected components of $M$).
\end{recall}

\begin{node}
We also recall we have the mod-2 ``fundamental class'' $[M]_{2}\in H_{n}(M;\ZZ_{2})$
which is just $[M]_{2}=(1,1,\dots,1)$.
\end{node}

\begin{node}
We have the ``evaluation map''
\begin{align}
\eval\colon & H^{n}(M;\ZZ_{2})\to\ZZ_{2}\\
&\alpha\mapsto\langle\alpha,[M]_{2}\rangle\nonumber.
\end{align}
\end{node}

\begin{puzzle}
How do we turn $w_{k}(M)$ into some [hopefully useful] number?
\end{puzzle}

\begin{definition}
Let $M$ be a smooth $n$-manifold.
Let $r_{1}$, \dots, $r_{n}\geq0$ be non-negative integers such that
$\sum^{n}_{k=1}kr_{k}=n$ (i.e.,
$r_{1}+2r_{2}+3r_{3}+\dots+nr_{n}=n$).
Observe that $w_{1}^{r_{1}}(\T M)w_{2}^{r_{2}}(\T M)(\dots)w_{n}^{r_{n}}(\T M)\in H^{n}(M;\ZZ_{2})$.
We define the \define{Steifel--Whitney Number} associated with
$(r_{1},\dots,r_{n})$ to be equal to
$\langle w_{1}^{r_{1}}(\T M)w_{2}^{r_{2}}(\T M)(\dots)w_{n}^{r_{n}}(\T M),[M]_{2}\rangle\in\ZZ_{2}$.
\end{definition}

\begin{remark}
Recall in Lecture~1, $\langle w_{n}(\T M),[M]_{2}\rangle=\chi(M)\bmod{2}$
is the Euler characteristic modulo~2.
\end{remark}

\begin{definition}
Let $M$, $M'$ be two smooth $n$-manifolds.
We say $M$ and $M'$ \define{Have the Same Stiefel--Whitney Numbers}
if
$\langle w_{1}^{r_{1}}(\T M)w_{2}^{r_{2}}(\T M)(\dots)w_{n}^{r_{n}}(\T M),[M]_{2}\rangle=\langle w_{1}^{r_{1}}(\T M')w_{2}^{r_{2}}(\T M')(\dots)w_{n}^{r_{n}}(\T M'),[M']_{2}\rangle$.
\end{definition}

\begin{example}
If $M$ is parallelizable, then $w_{0}(\T M)=1$ and $w_{i}(\T M)=0$ for
all $i\geq1$, which requires $r_{i}=0$ for all $i\geq1$.
\end{example}